\section{Apéndice: índice de términos y ruta de repaso}

\subsection{Índice de términos}

\begin{longtable}{p{0.23\textwidth}p{0.32\textwidth}p{0.36\textwidth}}
\toprule
Término & Explicación en una frase & Función en este episodio \\
\midrule
Agent & Entidad con límites que actúa en un entorno orientada a objetivos & Objeto central de todo el texto; abarca personas, animales, agentes de software y agentes del mundo digital. \\
Logical Agent & Agente inteligente basado en reglas lógicas y razonamiento simbólico & Muestra que Agent no es un concepto nuevo, sino un objetivo temprano de la IA. \\
Semantic Parsing & Convertir lenguaje natural en representaciones ejecutables como SQL, API o expresiones lógicas & Antecedente directo de Language Agent. \\
Language Agent & Agent que usa un LLM como interfaz de lenguaje y puede planificar y usar herramientas & Base técnica del OpenClaw Moment. \\
ReAct & Patrón que alterna Reasoning y Acting & Lleva al LLM de la respuesta única a un ciclo de interacción con el entorno. \\
Tool Use & El modelo invoca herramientas, API o funciones externas & Permite al modelo superar el límite de la respuesta de solo texto. \\
Computer Use Agent & Agent capaz de operar navegadores, escritorios, dispositivos móviles o entornos de código & Puerta de entrada importante hacia el universal digital agent. \\
World Model & Modelo interno del estado del entorno, de su causalidad y de las acciones posibles & Base del continual learning y de la confiabilidad. \\
Continual Learning & Actualizar las capacidades de forma continua a partir de nuevas interacciones & Candidato a eje principal del desarrollo de los Agent en 2026. \\
Expert Agent & Agent especializado que acumula habilidades en un dominio & Forma resultante que resuelve los problemas de confiabilidad, velocidad y costo. \\
\bottomrule
\end{longtable}

\subsection{Ruta de repaso}

Si solo se quiere repasar rápidamente este episodio, se puede leer en cuatro pasos: primero, la sección 2, para dominar la definición mínima de Agent; después, las secciones 3--4, para entender la historia técnica y la pila de Language Agent; luego, las secciones 5--7, para entender OpenClaw, la disolución de fronteras y el continual learning; por último, las secciones 8--10, para entender las bets de las grandes empresas tecnológicas, la evaluación en producción y la difusión social. Así, la pregunta «¿por qué los Agent se volvieron importantes de repente?» deja de ser una noticia de productos y se replantea como un problema de sistemas técnicos.

\begin{knowledgebox}{Conexión entre este episodio y el siguiente}
EP139 presenta el marco técnico de los Agent; la entrevista con Luo Fuli (罗福莉) en EP138 sitúa ese marco dentro de los problemas de un laboratorio de modelos: el postentrenamiento, la RL infra, la asignación de GPU y la igualdad dentro de la organización. Leer primero EP139 y después EP138 facilita entender por qué «el paradigma de los Agent depende mucho del postentrenamiento».
\end{knowledgebox}
